El juliol del 2020 la NASA va posar en marxa una missió espacial que, entre altres tasques, havia de fer arribar el vehicle d'exploració \textit{Perseverance} a la superfície de Mart. El 18 de febrer de 2021 va tenir lloc l'aterratge del vehicle. En la darrera etapa d'aquest procés complex d'aterratge, una grua fa baixar d'una manera controlada el vehicle des d'una altura de 20,0 m per sobre de la superfície de Mart. Durant tot aquest recorregut, la intensitat del camp gravitatori es pot considerar uniforme.

\begin{apartat}{1,25}
El valor absolut de la diferència de potencial gravitatori entre la superfície del planeta i un punt elevat 20,0 m per sobre de la superfície és 74,4 J/kg. A partir de la diferència de potencial, determineu el mòdul de la intensitat del camp gravitatori a la superfície de Mart. Quin dels esquemes següents ($a$, $b$, $c$ o $d$) representa les línies equipotencials a la superfície de Mart? Situeu en el diagrama triat la línia de menor potencial ($V_\mathrm{baix}$) i la de major potencial ($V_\mathrm{alt}$). Justifiqueu totes les respostes.

\figura[0.9]{fig1.png}
\end{apartat}

\begin{apartat}{1,25}
Si durant aquesta darrera etapa, el vehicle fa un descens de 20,0 m a una velocitat constant, quin treball ha fet la grua? Podeu negligir el treball fet per les forces de fregament.
\end{apartat}

\dades{%
  La massa del vehicle és 1\,025 kg.}
